% report_naver_v2.tex -- Workshop paper, English only (Korean to be added later).
% Compile with xelatex (twice, for cross-references).
%
% Final specification agreed after supervisor feedback (October 2026):
%   M1 Revisit  — binary logistic regression (6 vars)
%   M2 Rating   — multinomial logistic regression (5 vars, 3 categories)
% The ordered-logit M2 of earlier drafts was dropped because the Brant test
% rejected proportional odds for sentiment and ln(review_age).  Multinomial
% logit removes that constraint by estimating category-specific coefficients.
%
% Tables with static descriptives come from tables_naver.tex (generated by
% 10_01_descriptive_tables.py).  Model-result tables are typeset here from
% 09_FINAL_spec_diagnostics.py outputs under data/analysis/naver/workshop/FINAL/.

\documentclass[10pt,a4paper]{article}

\usepackage{fontspec}
\usepackage{xeCJK}
\setCJKmainfont{Malgun Gothic}
\setCJKsansfont{Malgun Gothic}
\xeCJKsetup{CJKspace=true}

\usepackage[top=2.0cm, bottom=2.2cm, left=2.2cm, right=2.2cm]{geometry}
\usepackage{booktabs}
\usepackage{amsmath}
\usepackage{enumitem}
\usepackage{microtype}
\usepackage{caption}
\usepackage{float}
\usepackage{longtable}
\usepackage{fancyhdr}
\usepackage{graphicx}
\usepackage[hidelinks]{hyperref}
\captionsetup[figure]{name=Figure, labelsep=period, font=small, labelfont=bf,
                     skip=4pt, justification=centering}

\captionsetup[table]{name=Table, labelsep=period, font=small, labelfont=bf,
                     skip=4pt, justification=centering}

% ─── Conference running header (match abstract template) ────────────────
\pagestyle{fancy}
\fancyhf{}
\renewcommand{\headrulewidth}{0.4pt}
\fancyhead[L]{\small 2026년 한국생산운영관리학회 추계학술대회 발표 초록집}
\fancyfoot[C]{\small -- \thepage{} --}

\linespread{1.04}
\setlength{\parindent}{0pt}
\setlength{\parskip}{5pt}

\newcommand{\sect}[1]{\vspace{10pt}\noindent{\Large\bfseries #1}\par\vspace{3pt}}
\newcommand{\subsect}[1]{\vspace{6pt}\noindent{\large\bfseries #1}\par\vspace{2pt}}
\newcommand{\todo}[1]{\par\noindent\fbox{\parbox{0.96\linewidth}{\small\textit{[SOURCE NEEDED: #1]}}}\par}

\title{}
\date{}

\begin{document}

% ═══════════════════════════════════════════════════════════════════════
% TITLE + AUTHOR (match the submitted abstract template)
% ═══════════════════════════════════════════════════════════════════════
\begin{center}
{\Large\textbf{An Empirical Analysis of Visit Context and\\[2pt]
Review Outcomes in Online Reviews of Korean Tourist Attractions}}

\vspace{12pt}

{\normalsize *Jorge Adan Soriano Hernandez (Department of Data Science,
Hoseo University, Graduate Student, adan\_25@naver.com)}

\end{center}

\vspace{10pt}

% ═══════════════════════════════════════════════════════════════════════
% ABSTRACT  — same text as the submitted abstract (no numeric results)
% Only methodological wording updated from "ordered" to "multinomial"
% to reflect the final specification estimated in the paper body.
% ═══════════════════════════════════════════════════════════════════════
\noindent\textbf{\large Abstract}

\vspace{4pt}

\noindent
This study aims to examine how platform-structured visit context
(companion type, visit purpose, whether the reviewer waited in line, and
whether a reservation was made) is reflected in the outcomes of online
reviews of Korean tourist attractions on Naver Place. Prior work has
focused on review text, sentiment, or topic structure, and the use of
platform-provided visit-context metadata as a joint predictor set to
explain review-level outcomes remains limited. Moreover, most existing
studies are restricted to a single attraction or category, leaving a gap
in review-level comparative analysis across heterogeneous destinations.
This study addresses two research questions. First, how do visit context
and reviewer characteristics shape the reviewer's declared satisfaction,
expressed as the star rating? Second, how do visit context and reviewer
characteristics relate to the revisit indicator? Both outcomes are
interpreted as different observable traces of an unobservable
satisfaction construct, with the star rating as a declared signal and
the revisit flag as a behavioural one grounded in prior experience. The
dataset consists of 14{,}401 Korean-language reviews from 29 major
Korean tourist attractions in and around Seoul and Busan, selected via
stratified sampling based on visitor volume. A multinomial logistic
regression of the star rating (grouped into TOP $=5$\textbigstar,
MID $4$--$4.5$\textbigstar, LOW $\le 3.5$\textbigstar) and a binary
logistic regression of the revisit indicator are employed, both using
visit context, reviewer characteristics, and place controls as common
predictors. This study contributes by proposing an analytic framework
that uses platform-structured visit context as a joint predictor set at
the review level, explaining both declared and behavioural traces of
satisfaction simultaneously.

\vspace{6pt}

\noindent\textbf{주제어 (Keywords):} review, visit context, revisit,
satisfaction, logistic regression

\sect{I. Introduction}

\subsect{1.1 Research Background}
Online reviews have become a decisive input to tourists' choice of where
to visit and whether to return. Before travelling, consumers routinely
consult reviews, photos and ratings on platforms such as TripAdvisor,
Google Maps, and in the Korean market, Naver Place and Kakao Map. These
platforms differ from each other in which data they expose: whereas most
Western platforms publish review text and a numeric rating, Naver Place
additionally records structured metadata about the visit itself --- who
the reviewer went with, why they came, whether they had to wait at the
entrance, and whether they had made a reservation. These platform-
captured fields turn each review into a small, machine-readable case
file of a visit.

South Korea's inbound tourism has grown substantially over the past
decade. Even against the backdrop of post-pandemic recovery, Korea
remains one of Asia's most-visited destinations, drawing visitors
across heterogeneous experiential segments. Research attributes the
growth partly to the Korean Wave (Hallyu) and the global spread of
Korean cultural content (Bae, Chang, Park \& Kim, 2017), and partly to
specific experiential segments such as food tourism (Chang, 2021; Ko,
Kang, Kang \& Lee, 2018). Naver Place, the dominant local review
platform for domestic Korean reviewers, accumulates several million
reviews across tourist attractions each year, offering an unusually
rich source for review-level analysis.

Within this context, understanding which contextual factors shape
\emph{how} visitors evaluate attractions and \emph{whether} they return
is of practical interest for destination managers and of academic
interest for consumer-behaviour researchers. Most existing work either
analyses review text in isolation --- treating the review as a
disembodied document --- or aggregates reviews up to a business level,
losing the review as a unit of analysis. The structured visit-context
metadata that Naver captures makes it possible, for the first time, to
model each review while keeping its visit context intact.

Throughout this paper, we refer to \textbf{visit context} as the set of
platform-captured dimensions of a visit: companion type (alone, partner,
friends, children, family, colleagues), visit purpose (travel, outing,
daily life, date, socializing), wait experience (binary), and
reservation (binary). These dimensions are platform-structured
metadata that few other review ecosystems record, and they constitute
the central predictor set examined throughout this study. We complement
them with reviewer-side variables (follower count, number of prior
reviews, number of photos uploaded with the review) and a text-based
sentiment score, so that each review is represented by a compact set of
visit, reviewer and text features.

\subsect{1.2 Research Questions}
We address two research questions:
\begin{enumerate}[leftmargin=16pt, itemsep=2pt]
  \item How do visit context and reviewer characteristics shape the
        probability that a review is written by a returning visitor?
  \item How do visit context and reviewer characteristics shape the
        probability of each star-rating category (five stars, four to
        four-and-a-half stars, three-and-a-half stars or fewer)?
\end{enumerate}

The two outcomes are interpreted as distinct observable traces of a
latent satisfaction construct: the star rating captures a reviewer's
\emph{declared} satisfaction at the time of writing, while the revisit
flag captures a \emph{behavioural} trace of satisfaction accumulated
across prior visits. We model them in parallel with a near-identical
predictor set so that differences in coefficients across the two models
reflect differences in the outcomes themselves, not differences in model
specification.

\subsect{1.3 Contribution}
This paper makes three contributions. \emph{(i)~Platform-structured
visit context as a joint predictor set.} We exploit Naver Place's
structured metadata (companion, purpose, wait, reservation) as review-
level predictors, rather than restricting the analysis to review text
alone. \emph{(ii)~Parallel modelling of two outcomes.} We estimate two
regressions with a shared core of predictors, allowing direct comparison
of how each contextual factor shapes revisit versus star rating.
\emph{(iii)~Multinomial logit for the three-class rating.} Because an
ordered-logit specification of the rating failed the Brant test for
proportional odds (see \S3.5), we adopt a multinomial logit that permits
category-specific coefficients, revealing asymmetric effects of key
predictors across rating categories.

\sect{II. Theoretical Background}

\subsect{2.1 Visit Context and Review Behaviour}
Lee, Jeong and Kim (2026) show that review-side characteristics shape
business-level performance in Korean platform data, using Naver Smart
Place at the business-week level ($N = 2{,}662{,}298$ business-weeks).
Their analysis establishes that information diversity (the uniqueness
of keyword content across reviews) and review sentiment both predict
platform traffic (page views and total clicks), and documents an
important interaction with Naver's March 2023 keyword-expansion update:
the effect of text length on traffic reverses sign after the update,
and sentiment moves from negative to positive. Their contribution is
to show that \emph{what} reviews contain matters above and beyond
\emph{how many} reviews exist. However, their aggregation to business-
week level means the review itself --- and the visit context in which
it was written --- are not directly observed. The present paper takes
their insight to the review level, keeping the review as the unit of
analysis and asking what review- and visit-level predictors drive two
behavioural outcomes: return visits and star ratings.

Tourism-segmentation research complements this by treating visit
purpose and companion type as the two dimensions that most strongly
shape what visitors experience and how they later evaluate it. Ko,
Kang, Kang and Lee (2018) use factor-cluster segmentation on visit
purpose to classify foreign tourists' perceptions of Korean food tours;
Chang (2021) treats food tourism itself as a distinct experiential
segment with its own antecedents and outcomes; Bae, Chang, Park and
Kim (2017) connect tourism-segment growth to the Korean Wave. These
studies motivate the inclusion of visit purpose (``tourist'' indicator)
and companion type (``alone'' indicator) as review-level predictors in
the present analysis.

Wait experience and reservation, by contrast, operate as operational
signals of the service encounter and are prominent in the service-
quality literature (SERVQUAL; Parasuraman, Zeithaml \& Berry, 1988),
which identifies five dimensions of perceived service quality:
tangibles, reliability, responsiveness, assurance, and empathy. Waiting
at the entrance is a direct violation of responsiveness, and prior work
on tourism and hospitality shows that queueing has strong, often
disproportionate effects on satisfaction and behavioural intention.
Reservation, in contrast, is a signal of planning and commitment; it
also raises baseline expectations about service.

\subsect{2.2 Satisfaction as a Latent Construct}
Customer satisfaction is routinely conceptualised in the consumer-
behaviour literature as a \emph{latent} construct with multiple
observable indicators. The classic measurement perspective
(e.g., Oliver's expectancy-disconfirmation model) treats satisfaction
as an unobservable psychological state inferred from several behavioural
and declarative traces: a star rating, a textual sentiment, a repurchase
intention, a repeat visit. No single trace fully reveals the latent
construct, and the traces are typically only moderately correlated.

In the review-platform setting, two particularly useful observable
indicators coexist in the data: the star rating (a \emph{declared}
indicator produced at the moment of writing) and the revisit flag
(a \emph{behavioural} indicator that aggregates past visits). These
two indicators are related but need not move together: a reviewer may
award five stars yet never return (for example, a tourist on a one-time
trip who found the attraction worth their time but has no reason to
return), or they may return repeatedly while rating the attraction with
less than the maximum (for example, a local who finds the venue useful
without being enthusiastic). This paper's design is to model both
indicators in parallel with a near-identical predictor set, so that
differences in coefficients across the two models can be attributed to
differences in what each outcome measures rather than to differences in
model specification.

The weak empirical correlation between the two outcomes in our sample
($r = 0.08$, see \S3.4) supports this framing: although both revisit
and rating plausibly reflect satisfaction, they capture distinct
aspects of it and share less than 1\% of their variance. Divergent
patterns of coefficients across the two models are therefore substantive
findings about how different contextual factors route through to the
two observable traces, not artefacts of modelling choice.

\subsect{2.3 Sentiment Classification and Review Helpfulness}
Review-text sentiment is scored with a BERT-family Korean sentiment
classifier (\texttt{hun3359/klue-bert-base-sentiment}), a model
fine-tuned on Korean consumer-review data. BERT-based sentiment models
substantially outperform earlier bag-of-words or rule-based approaches
in consumer-review settings, reaching 91\% accuracy in product-
recommendation contexts (Karabila, Darraz, EL-Ansari, Alami \& EL
Mallahi, 2024). We treat sentiment as a text-based predictor and
include it in both outcome models.

For the rating outcome, we acknowledge upfront that the association
between sentiment and rating is partly mechanical: both are produced by
the same reviewer in the same session, so a review that reads positively
and a high star rating share a common authorship. The sentiment
coefficient in the rating model should therefore be read as a text-
level control, not as a causal effect. For the revisit outcome, the
mechanical overlap is weaker --- the revisit flag reflects behaviour
across multiple visits, not just the current one --- so the sentiment
coefficient there is more interpretable as a review-content effect.

\subsect{2.4 Research Gap}
Existing studies either analyse review text in isolation or aggregate
reviews to business-week performance; few connect platform-captured
visit context to review-level outcomes. Furthermore, when the star
rating is modelled as an ordinal outcome, the proportional-odds
assumption underpinning the ordered-logit specification is rarely
tested, which can mask important category-specific effects. The present
paper addresses both gaps: it uses the Naver-recorded visit context as
a predictor set at the review level, and models the rating as a
multinomial outcome after an explicit test of proportional odds (see
\S3.5). This combination lets us compare two parallel outcomes using
mostly the same predictors while permitting the rating model to recover
heterogeneous effects across the three rating categories.

\sect{III. Methodology}

\subsect{3.1 Data}
We analyse 14{,}401 Korean-language reviews from 29 major Korean tourist
attractions in and around Seoul and Busan, collected by scraping Naver
Place for each attraction individually. Of these, 11{,}165 reviews carry
free-text content and 4{,}044 carry both text and a star rating; the
latter is the sample for the rating model (M2). Attractions were
selected via stratified sampling across typological categories (palaces,
traditional quarters, markets, museums, modern landmarks, beaches,
parks, theme parks), with a minimum of 30 Korean-language reviews per
attraction.

\subsect{3.2 Variables}

Each review contributes one observation. The dependent variables are:
\textbf{Revisit}, a binary indicator equal to~1 if the reviewer reports
having visited the attraction at least twice; and \textbf{Rating},
grouped from the original 0.5--5.0 star scale into three categories to
support the multinomial specification:
\begin{itemize}[leftmargin=14pt, itemsep=0pt]
  \item \textbf{LOW}: rating $\le 3.5$\textbigstar\ (276 reviews, 6.8\%)
  \item \textbf{MID}: rating $\in \{4.0, 4.5\}$\textbigstar\ (728
        reviews, 18.0\%)
  \item \textbf{TOP}: rating $= 5.0$\textbigstar\ (3{,}040 reviews, 75.2\%)
\end{itemize}

Predictors are drawn from three groups of platform-captured fields:
\emph{visit context} (reviewer-reported fields about the visit),
\emph{reviewer characteristics} (platform-level attributes of the
reviewer's profile), and a derived text-based predictor
(\emph{sentiment}). Table~\ref{tab:vars} lists the predictors used
in each model.

\begin{table}[H]\centering
\caption{Predictor variables by model. Three predictors (sentiment,
waited, ln(followers)) are shared across both models; the remaining
variables are specific to one.}
\label{tab:vars}
\footnotesize
\begin{tabular}{@{}p{3.4cm}p{3.6cm}p{3.4cm}cc@{}}
\toprule
Variable & Description & Scale & M1 & M2 \\
\midrule
sentiment            & BERT sentiment of the text         & continuous $[-1, 1]$   & \checkmark & \checkmark \\
waited               & Had to queue at entrance           & binary $\{0, 1\}$        & \checkmark & \checkmark \\
ln(1+followers)      & Reviewer's follower count           & continuous               & \checkmark & \checkmark \\
\midrule
is\_tourist          & Visit purpose was ``travel''        & binary $\{0, 1\}$        & \checkmark & -- \\
comp\_alone          & Visit was solo                      & binary $\{0, 1\}$        & \checkmark & -- \\
ln(1+reviews)        & Reviewer's prior review count       & continuous               & \checkmark & -- \\
\midrule
ln(1+photos)         & Photos attached to the review       & continuous               & -- & \checkmark \\
season\_autumn       & Visit was in autumn                 & binary $\{0, 1\}$        & -- & \checkmark \\
\bottomrule
\end{tabular}
\end{table}

Variable selection followed an exhaustive sweep comparing more than
eighty specifications across both outcomes (details in supplementary
material). The final specification retains, for each outcome, the
predictors that are significant and whose inclusion improves AIC. The
three shared predictors emerged as the only variables significant across
both outcomes under every subset of the candidate pool.

\subsect{3.3 Descriptive Statistics}
Table~\ref{tab:desc} summarises the two outcome variables, the three
rating categories, and the eight predictor variables used across the
two models.

\begin{table}[H]\centering
\caption{Descriptive statistics for outcomes and predictors.
``N (M1)'' is the sample for the revisit model (reviews with text);
``N (M2)'' is the sample for the rating model (reviews with text
\emph{and} a star rating).}
\label{tab:desc}
\footnotesize
\begin{tabular}{@{}lrrrrrr@{}}
\toprule
Variable             & N (M1)   & Mean   & SD    & Min    & Max   & N (M2) \\
\midrule
\multicolumn{7}{@{}l}{\textit{Outcomes}} \\
Revisit (0/1)                 & 11,165 & 0.226 & 0.418 & 0 & 1 & --- \\
Five-star (0/1)               & ---    & 0.752 & 0.432 & 0 & 1 & 4,044 \\
Rating category: LOW ($\le 3.5$\textbigstar) & ---    & 0.068 & ---   & 0 & 1 & 276 \\
Rating category: MID (4--4.5\textbigstar) & ---    & 0.180 & ---   & 0 & 1 & 728 \\
Rating category: TOP (5\textbigstar) & ---    & 0.752 & ---   & 0 & 1 & 3,040 \\
\midrule
\multicolumn{7}{@{}l}{\textit{Predictors --- visit context}} \\
Waited (0/1)                  & 11,165 & 0.079 & 0.270 & 0 & 1 & 4,044 \\
Tourist (0/1)                 & 11,165 & 0.163 & 0.369 & 0 & 1 & 4,044 \\
Alone (0/1)                   & 11,165 & 0.110 & 0.312 & 0 & 1 & 4,044 \\
Season: autumn (0/1)          & 11,165 & 0.301 & 0.459 & 0 & 1 & 4,044 \\
\midrule
\multicolumn{7}{@{}l}{\textit{Predictors --- reviewer characteristics}} \\
ln(1 + reviews)               & 11,165 & 5.386 & 1.927 & 0.69 & 9.86 & 4,044 \\
ln(1 + followers)             & 11,165 & 1.895 & 2.245 & 0.00 & 9.17 & 4,044 \\
ln(1 + photos)                & 11,165 & 0.100 & 0.375 & 0.00 & 3.05 & 4,044 \\
\midrule
\multicolumn{7}{@{}l}{\textit{Predictor --- text signal}} \\
Sentiment                     & 11,165 & 0.653 & 0.457 & $-0.97$ & 0.99 & 4,044 \\
\bottomrule
\end{tabular}
\end{table}

Approximately 23\% of reviews with text are from returning visitors;
75\% of rated reviews award the maximum of five stars. Only 8\% of
reviews report having waited at the entrance; 16\% are tourist-purpose
visits and 11\% are solo visits. Autumn accounts for 30\% of visits in
the sample. Sentiment is skewed positive (mean $0.65$ on a $[-1, 1]$
scale), consistent with the general positivity of online tourism
reviews.

\subsect{3.4 Correlation Analysis and Multicollinearity}
Table~\ref{tab:corr} reports the pairwise correlation matrix for the
eight predictors in the rating-model sample ($N = 4{,}044$).
Table~\ref{tab:vif} reports the variance inflation factors (VIFs).

\begin{table}[H]\centering
\caption{Pairwise correlations among the eight predictor variables
(M2 sample, $N = 4{,}044$). All correlations are below 0.40 in absolute
value except for ln(reviews) with ln(followers) ($r = 0.73$).}
\label{tab:corr}
\scriptsize
\setlength{\tabcolsep}{3pt}
\begin{tabular}{@{}lrrrrrrrr@{}}
\toprule
                       & Wait  & Tour  & Sent   & Alone & lnRev & lnFol & lnPho & Autum \\
\midrule
Waited (Wait)          &  1.00 &       &        &       &       &       &       &       \\
Tourist (Tour)         &  0.09 &  1.00 &        &       &       &       &       &       \\
Sentiment (Sent)       & -0.08 & -0.06 &  1.00  &       &       &       &       &       \\
Alone                  & -0.06 &  0.09 & -0.09  &  1.00 &       &       &       &       \\
ln(reviews) (lnRev)    & -0.08 &  0.00 & -0.09  &  0.20 &  1.00 &       &       &       \\
ln(followers) (lnFol)  & -0.10 & -0.02 & -0.12  &  0.28 &  0.73 &  1.00 &       &       \\
ln(photos) (lnPho)     &  0.01 &  0.01 &  0.03  &  0.01 &  0.08 &  0.06 &  1.00 &       \\
Season autumn (Autum)  & -0.03 & -0.04 &  0.02  & -0.05 & -0.10 & -0.07 & -0.03 &  1.00 \\
\bottomrule
\end{tabular}
\end{table}

\begin{table}[H]\centering
\caption{Variance inflation factors (M2 sample). All VIFs are below
2.5, well under the common threshold of 5--10 for multicollinearity
concern.}
\label{tab:vif}
\footnotesize
\begin{tabular}{@{}lr lr@{}}
\toprule
Variable          & VIF  & Variable          & VIF \\
\midrule
Waited             & 1.03 & ln(1+reviews)     & 2.17 \\
Tourist            & 1.02 & ln(1+followers)   & 2.27 \\
Sentiment          & 1.03 & ln(1+photos)      & 1.01 \\
Alone              & 1.10 & Season: autumn     & 1.01 \\
\bottomrule
\end{tabular}
\end{table}

The highest correlation, between ln(reviews) and ln(followers)
($r = 0.73$), is intuitive: more-active reviewers also tend to
accumulate larger follower bases. Despite this correlation, both
variables are retained in M1 only (not jointly in M2), and the
individual VIFs ($2.17$ and $2.27$) remain comfortably below standard
thresholds, so neither variable's standard error is unduly inflated.

Critically for the paper's framing, the two outcomes themselves are
only weakly correlated: $r = 0.08$ between Revisit and the five-star
indicator in the shared sample ($N = 4{,}044$, $p < 0.001$). This low
correlation is the empirical justification for modelling the two
outcomes separately: although both are plausibly indicators of a latent
satisfaction construct, they capture distinct aspects of it and share
only about 0.6\% of their variance.

\subsect{3.5 Models}
Both outcome models include place fixed effects $\delta_j$ (28 dummies,
Gyeongbokgung as reference) and standard errors clustered by reviewer
(\texttt{user\_id}). Let $i$ index review and $j$ index place; let
$\Lambda(z) = 1/(1+e^{-z})$ denote the logistic function.

\textbf{Model 1 --- Binary Logit (Revisit).} The outcome
$\text{Revisit}_{ij}$ equals $1$ if the reviewer's reported visit count
is $\ge 2$ and $0$ otherwise.

\begin{align*}
\operatorname{logit} P(\text{Revisit}_{ij} = 1)
   = \; &\alpha
       + \beta_{1}\,\text{waited}_{ij}
       + \beta_{2}\,\text{is\_tourist}_{ij}
       + \beta_{3}\,\text{sentiment}_{ij} \\
    &+ \beta_{4}\,\text{comp\_alone}_{ij}
       + \beta_{5}\,\ln(1 + \text{reviews}_{ij})
       + \beta_{6}\,\ln(1 + \text{followers}_{ij})
       + \delta_{j}
\end{align*}

\textbf{Model 2 --- Multinomial Logit (Rating).} Let
$Y_{ij} \in \{\text{LOW}, \text{MID}, \text{TOP}\}$ denote the
three-category rating, with TOP as the reference category. For each
non-reference category $k \in \{\text{MID}, \text{LOW}\}$:

\begin{align*}
\log\!\frac{P(Y_{ij} = k)}{P(Y_{ij} = \text{TOP})}
   = \; &\alpha_{k}
       + \gamma_{k1}\,\text{sentiment}_{ij}
       + \gamma_{k2}\,\text{waited}_{ij}
       + \gamma_{k3}\,\ln(1 + \text{followers}_{ij}) \\
    &+ \gamma_{k4}\,\ln(1 + \text{photos}_{ij})
       + \gamma_{k5}\,\text{season\_autumn}_{ij}
       + \delta_{jk}
\end{align*}

Each non-reference category is associated with its own intercept
$\alpha_{k}$, its own vector of slope coefficients
$(\gamma_{k1}, \ldots, \gamma_{k5})$, and its own set of place
fixed-effects $\delta_{jk}$. This specification yields two parallel
sets of coefficients (MID vs TOP, LOW vs TOP) that are allowed to
differ freely, rather than being constrained to a single proportional
form.

\textbf{Why multinomial and not ordered logit?} We initially specified M2
as an ordered logistic regression. The Brant test rejected the
proportional-odds assumption for the full model ($\chi^2 = 18.8$, df$=7$,
$p = 0.009$), with sentiment and ln(review\_age) violating individually
($p < 0.01$). Multinomial logit removes this assumption by estimating
separate coefficients for each category comparison; the trade-off is
that the ordinal structure of the rating is not imposed on the fit.

\sect{IV. Results}

\subsect{4.0 Model Fit Summary}

Table~\ref{tab:fit} summarises the fit of the two models side by side.
Both models reach the ``good'' range for discrete-choice models in
consumer research (pseudo-$R^2$ above 0.10, with 0.20--0.40 considered
excellent per McFadden). Out-of-sample metrics come from five-fold
cross-validation grouped by reviewer (no reviewer appears in both
training and test folds).

\begin{table}[H]\centering
\caption{Model fit summary. Pseudo-$R^2$ is McFadden's; AIC and BIC are
Akaike and Bayesian information criteria (lower is better). AUC is the
area under the ROC curve from cross-validation; within-$\pm 1$ accuracy
is the proportion of M2 predictions within one category of the actual
category.}
\label{tab:fit}
\footnotesize
\begin{tabular}{@{}l l l@{}}
\toprule
Statistic                     & M1 Revisit         & M2 Rating \\
\midrule
Model type                     & Binary logit       & Multinomial logit \\
Outcome                        & Revisit (0/1)      & Rating (LOW / MID / TOP) \\
Reference category             & no revisit (0)     & TOP (5\textbigstar) \\
N observations                 & 11,165             & 4,044 \\
Predictors                     & 6                  & 5 (per category) \\
Place fixed effects            & 28                 & 28 (per category) \\
pseudo-$R^2$                   & 0.201              & 0.190 \\
AIC                            & 9{,}601.5          & 4{,}753.6 \\
BIC                            & 9{,}857.7          & 5{,}157.1 \\
Log-likelihood                 & $-4{,}765.7$       & $-2{,}312.8$ \\
AUC / Exact accuracy (CV)      & 0.80               & 0.77 \\
Within-$\pm 1$ accuracy (M2)   & ---                & 0.95 \\
\bottomrule
\end{tabular}
\end{table}

\subsect{4.1 Model 1 — Revisit}

Table~\ref{tab:m1} reports the coefficient estimates for Model 1.
All six predictors are significant at $p < 0.01$. Pseudo-$R^2$ is
0.201 and the out-of-sample AUC from five-fold cross-validation grouped
by reviewer is 0.78.

\begin{table}[H]\centering
\caption{Model 1: Binary logistic regression of Revisit.
$N = 11{,}165$; pseudo-$R^2 = 0.201$; AIC $= 9{,}602$; BIC $= 9{,}858$.
Place fixed effects included; standard errors clustered by reviewer.}
\label{tab:m1}
\footnotesize
\begin{tabular}{@{}lrrr@{}}
\toprule
Variable & OR & 95\% CI & $p$ \\
\midrule
Waited to enter               & 0.54 & [0.37, 0.79] & 0.001 ** \\
Tourist (purpose = travel)    & 0.55 & [0.42, 0.73] & $<\,.001$ *** \\
Sentiment ($-1$ to $+1$)      & 0.74 & [0.64, 0.87] & $<\,.001$ *** \\
Alone (solo visit)             & 2.26 & [1.67, 3.07] & $<\,.001$ *** \\
ln(1 + reviewer's reviews)    & 1.31 & [1.17, 1.48] & $<\,.001$ *** \\
ln(1 + reviewer's followers)  & 1.23 & [1.13, 1.35] & $<\,.001$ *** \\
\bottomrule
\end{tabular}
\end{table}

\begin{figure}[H]\centering
\includegraphics[width=0.85\linewidth]{../../data/analysis/naver/workshop/FINAL/figures/forest_m1.png}
\caption{Model 1 odds ratios with 95\% confidence intervals (log
scale). Blue markers indicate $p < 0.05$. Dashed line at OR $= 1$ is
the null-effect reference.}
\label{fig:forest_m1}
\end{figure}

The strongest negative effect is \textbf{waiting at the entrance}:
reviews from reviewers who report having waited are about half as likely
to come from returning visitors (OR $= 0.54$). This is the paper's
clearest managerial finding: queueing is the single operational variable
in this specification that is directly actionable by attraction
operators.

Two \textbf{reviewer-identity variables} enter positively:
ln(1$+$reviews) (OR $= 1.31$) and ln(1$+$followers) (OR $= 1.23$).
Reviewers with more platform history or larger networks are
disproportionately likely to be return visitors. This likely reflects
selection — committed reviewers are both more active online and more
likely to make repeat visits — but it also means that platform-level
reviewer identity is informative about return behaviour independently of
the specific visit.

\textbf{Visit-type effects} are large and in opposite directions.
Tourist-purpose visits are 45\% less likely to come from returning
visitors (OR $= 0.55$), consistent with the fact that tourists typically
visit each site once. Conversely, solo visits more than double the odds
of a return visit (OR $= 2.26$), suggesting that solo visitors are
either locals who use attractions more regularly or committed enthusiasts
who return independently.

Sentiment enters \emph{negatively} in M1 (OR $= 0.74$): more positive
sentiment in a review is associated with lower odds of the review being
from a returning visitor. The sign is modest but statistically clear.
One plausible reading is that first-time visitors, having nothing to
compare against, write more uniformly positive reviews, while repeat
visitors introduce more nuance into their text.

\textbf{Classification performance.} At the 0.5 probability threshold,
the model achieves 82\% accuracy against a 77\% baseline (always predict
non-revisit). Balanced accuracy is 0.65; precision for revisits is 0.68
and recall 0.35, reflecting the class imbalance (23\% revisits). The
AUC of 0.78 (five-fold cross-validation grouped by reviewer; see
Table~\ref{tab:cm1}) is in the ``good'' range for behavioural models in
consumer research.

\begin{table}[H]\centering
\caption{Model 1 confusion matrix at $p = 0.50$ threshold
($N = 11{,}165$).}
\label{tab:cm1}
\footnotesize
\begin{tabular}{@{}lrrr@{}}
\toprule
                     & Predicted = 0 & Predicted = 1 & Row total \\
\midrule
Actual = 0 (no revisit) & 8{,}226 &   417 & 8{,}643 \\
Actual = 1 (revisit)    & 1{,}643 &   879 & 2{,}522 \\
\midrule
Column total            & 9{,}869 & 1{,}296 & 11{,}165 \\
\bottomrule
\end{tabular}
\end{table}

\subsect{4.2 Model 2 — Rating}

Table~\ref{tab:m2} reports Model 2 coefficients as two separate
columns of odds ratios: MID versus TOP, and LOW versus TOP. Pseudo-$R^2$
is 0.190.

\begin{table}[H]\centering
\caption{Model 2: Multinomial logistic regression of Rating category.
$N = 4{,}044$; pseudo-$R^2 = 0.190$; AIC $= 4{,}754$; BIC $= 5{,}157$.
Odds ratios relative to the TOP (five-star) baseline. Place fixed
effects included; SE clustered by reviewer.}
\label{tab:m2}
\footnotesize
\begin{tabular}{@{}l rr rr@{}}
\toprule
                            & \multicolumn{2}{c}{MID vs TOP} & \multicolumn{2}{c}{LOW vs TOP} \\
\cmidrule(lr){2-3} \cmidrule(l){4-5}
Variable                    & OR    & $p$          & OR    & $p$ \\
\midrule
Sentiment ($-1$ to $+1$)    & 0.44  & $<\,.001$ ***  & 0.04  & $<\,.001$ *** \\
Waited to enter             & 1.21  & 0.326        & 3.10  & $<\,.001$ *** \\
ln(1 + reviewer's followers)& 0.95  & 0.287        & 0.71  & $<\,.001$ *** \\
ln(1 + photos)              & 0.81  & 0.197        & 0.44  & 0.043 *      \\
Season: autumn              & 0.39  & $<\,.001$ ***  & 0.38  & $<\,.001$ *** \\
\bottomrule
\end{tabular}
\end{table}

\begin{figure}[H]\centering
\includegraphics[width=1.0\linewidth]{../../data/analysis/naver/workshop/FINAL/figures/forest_m2.png}
\caption{Model 2 odds ratios with 95\% confidence intervals (log
scale). Left panel: MID vs TOP. Right panel: LOW vs TOP. The
asymmetric magnitudes across panels are the main empirical signature
of the multinomial specification.}
\label{fig:forest_m2}
\end{figure}

The multinomial decomposition reveals that predictors operate
\textbf{asymmetrically across the two category comparisons}.

\textbf{Sentiment} has the strongest effect in both comparisons, but
far stronger when distinguishing LOW from TOP (OR $= 0.04$) than
MID from TOP (OR $= 0.44$). A one-unit increase in sentiment
multiplies the odds of a LOW rating (relative to TOP) by 0.04 — that
is, satisfied-sounding text is roughly 25 times less likely to accompany
a low rating than a top rating. Because sentiment and the rating are
produced from the same reviewer in the same session, this association
is partly mechanical and should not be interpreted as a causal effect.
We retain sentiment as a text-level control.

\textbf{Waiting at the entrance} shows a striking asymmetric pattern.
It is not significantly associated with the MID-vs-TOP comparison
($p = 0.33$), but strongly increases the odds of a LOW rating relative
to TOP (OR $= 3.10$, $p < 0.001$). In practice, having to wait is a
\emph{disaster signal} rather than a mediocrity signal: waiting is a
strong predictor of a very low rating, but it does not systematically
push reviewers toward intermediate ratings. The managerial implication
mirrors M1: waiting has a specific and strong negative association,
here with the probability of a severely low rating.

\textbf{Reviewer influence} (ln(followers)) also shows an asymmetric
pattern: not significant for MID vs TOP (OR $= 0.95$), but significantly
protective against LOW ratings (OR $= 0.71$). More-followed reviewers
appear less likely to write LOW ratings, perhaps because of a tacit
sense of review-platform norms or reputational self-consciousness.

\textbf{Photos.} ln(photos) is a weak predictor of MID vs TOP
($p = 0.20$) but reaches significance for LOW vs TOP (OR $= 0.44$,
$p = 0.04$): reviews with photos are about half as likely to be in the
LOW category.

\textbf{Autumn effect.} Visits during the autumn season (September to
November) have significantly lower odds of being in either MID or LOW
category than in TOP, with similar magnitudes (OR $= 0.39$ and 0.38
respectively). This is consistent with Korean autumn — foliage and
pleasant weather — being perceived as the peak season for the
attraction categories in our sample.

\textbf{Classification performance.} Table~\ref{tab:cm2} shows the M2
confusion matrix at the modal-prediction rule. Exact-category accuracy
is 77\% versus a 75\% baseline (always predict TOP). Within-$\pm 1$
accuracy — the proportion of predictions that are at most one category
away from the actual category — is \textbf{95\%}, which is strong for
a three-class ordinal problem.

\begin{table}[H]\centering
\caption{Model 2 confusion matrix (predictions by modal class).}
\label{tab:cm2}
\footnotesize
\begin{tabular}{@{}l rrr r@{}}
\toprule
                              & Pred TOP & Pred MID & Pred LOW & Row total \\
\midrule
Actual TOP (5\textbigstar)      & 2{,}982 &   10 &   48 & 3{,}040 \\
Actual MID (4--4.5\textbigstar) &    684 &    9 &   35 &   728 \\
Actual LOW ($\le 3.5$\textbigstar)&   138 &    1 &  137 &   276 \\
\midrule
Column total                    & 3{,}804 &   20 &  220 & 4{,}044 \\
\bottomrule
\end{tabular}
\end{table}

The off-diagonal structure reveals an interesting pattern. The model
assigns almost no observations to the MID class, routing most reviews
either to TOP or LOW. Per-class recall is 0.98 for TOP, 0.50 for LOW,
and only 0.01 for MID. This is a well-known limitation of discrete-
choice classification on heavily imbalanced outcomes: the model treats
MID as the \emph{indeterminate zone} between TOP and LOW rather than a
clearly distinguishable category. The 95\% within-$\pm 1$ accuracy
confirms that although exact-category assignment struggles with MID,
most prediction errors are off by only one category.

\subsect{4.3 Interpreting the Two Models Jointly}

Comparing the two models side by side (Table~\ref{tab:joint}) reveals
both convergence and divergence in how predictors operate on the two
outcomes.

\begin{table}[H]\centering
\caption{Side-by-side comparison of significant effects (summary).
Direction notation: $\uparrow$ increases outcome probability,
$\downarrow$ decreases, --~not significant.}
\label{tab:joint}
\footnotesize
\begin{tabular}{@{}l cc@{}}
\toprule
Predictor                   & M1 (Revisit)  & M2 (Rating) \\
\midrule
Waited                      & $\downarrow\downarrow$ (OR 0.54) & $\downarrow$ in LOW (OR 3.10) \\
Sentiment                    & $\downarrow$ (OR 0.74)             & $\uparrow\uparrow\uparrow$ in TOP (OR 0.04 for LOW) \\
ln(followers)                & $\uparrow$ (OR 1.23)                & protective against LOW (OR 0.71) \\
Tourist                       & $\downarrow\downarrow$ (OR 0.55) & --- (not in M2) \\
Alone                         & $\uparrow\uparrow$ (OR 2.26)       & --- (not in M2) \\
ln(reviews)                   & $\uparrow\uparrow$ (OR 1.31)       & --- (not in M2) \\
ln(photos)                    & --- (not in M1)                    & protective against LOW (OR 0.44) \\
Season: autumn                & --- (not in M1)                    & $\downarrow\downarrow$ in LOW/MID (both OR $\approx 0.4$) \\
\bottomrule
\end{tabular}
\end{table}

The three shared predictors (waited, sentiment, ln(followers))
illustrate the \emph{divergent-trace} interpretation we proposed in
\S2.2. Waiting is strongly negative for revisit, and specifically a
disaster signal in rating. Sentiment is modestly negative for revisit
(nuanced-but-positive reviewers are more likely to be returning) but
massively positive for high ratings (satisfied-sounding text accompanies
high ratings almost mechanically). Follower count is positive for both,
pointing to reviewer identity rather than visit characteristics.

The low correlation between the two outcomes ($r = 0.08$, see \S3.4)
confirms that these shared predictors are not forcing an identical story
on both models. If revisit and rating were essentially the same
underlying variable, the models would recover the same coefficients up
to scale; instead, they recover different coefficient patterns reflecting
the distinct behavioural and declared traces of satisfaction we outlined
in \S2.2.

\subsect{4.4 Cross-validation Details}

Table~\ref{tab:cv} reports the per-fold cross-validation metrics for
both models. The out-of-sample AUC for M1 varies from 0.75 to 0.84
across folds (mean 0.78); the exact accuracy for M2 varies from 0.76
to 0.80 (mean 0.77). No single fold exhibits markedly degraded
performance, suggesting the models generalise consistently across
reviewer subsets.

\begin{table}[H]\centering
\caption{Cross-validation detail (5-fold grouped by reviewer; no
reviewer appears in both the training and test fold).}
\label{tab:cv}
\footnotesize
\begin{tabular}{@{}c rrrrrr | rr@{}}
\toprule
     & \multicolumn{6}{c|}{M1 Revisit} & \multicolumn{2}{c}{M2 Rating} \\
Fold & N test  & AUC   & Brier & Acc.  & Bal.~Acc. & F1    & Exact & $\pm 1$ \\
\midrule
1    & 2{,}233 & 0.839 & 0.136 & 0.777 & 0.574     & 0.281 & 0.755 & 0.953 \\
2    & 2{,}233 & 0.797 & 0.131 & 0.824 & 0.659     & 0.477 & 0.802 & 0.962 \\
3    & 2{,}233 & 0.751 & 0.153 & 0.781 & 0.591     & 0.335 & 0.759 & 0.942 \\
4    & 2{,}233 & 0.759 & 0.147 & 0.785 & 0.573     & 0.284 & 0.768 & 0.958 \\
5    & 2{,}233 & 0.756 & 0.145 & 0.806 & 0.629     & 0.417 & 0.757 & 0.954 \\
\midrule
Avg. & 11{,}165 & 0.780 & 0.142 & 0.795 & 0.605 & 0.359 & 0.768 & 0.954 \\
\bottomrule
\end{tabular}
\end{table}

\subsect{4.5 Robustness Checks}

Table~\ref{tab:rob} reports alternative specifications of both models,
toggling individual predictors in or out of the main specification.
These checks confirm that the main-model fit statistics are stable to
small changes in the predictor set.

\begin{table}[H]\centering
\caption{Robustness alternatives. Each row estimates one specification;
the first row of each panel is the main model reported in \S4.1--4.2.}
\label{tab:rob}
\footnotesize
\begin{tabular}{@{}l ccrr@{}}
\toprule
Specification                         & N vars & pseudo-$R^2$ & AIC     & BIC     \\
\midrule
\multicolumn{5}{@{}l}{\textit{M1 Revisit alternatives}} \\
Main (6 vars)                          & 6      & 0.2010       & 9{,}601.5 & 9{,}857.7 \\
+ ln(review\_age)                      & 7      & 0.2133       & 9{,}457.6 & 9{,}721.2 \\
+ season\_autumn                       & 7      & 0.2013       & 9{,}600.5 & 9{,}864.0 \\
$-$ comp\_alone                        & 5      & 0.1912       & 9{,}717.5 & 9{,}966.4 \\
$-$ sentiment                          & 5      & 0.1989       & 9{,}624.8 & 9{,}873.7 \\
\midrule
\multicolumn{5}{@{}l}{\textit{M2 Rating alternatives}} \\
Main (5 vars)                          & 5      & 0.1904       & 4{,}753.6 & 5{,}157.1 \\
+ ln(review\_age)                      & 6      & 0.1920       & 4{,}748.5 & 5{,}164.6 \\
+ is\_tourist                           & 6      & 0.1924       & 4{,}746.1 & 5{,}162.2 \\
+ comp\_alone                          & 6      & 0.1904       & 4{,}757.6 & 5{,}173.7 \\
$-$ season\_autumn                     & 4      & 0.1749       & 4{,}838.5 & 5{,}229.4 \\
$-$ ln\_photos                         & 4      & 0.1885       & 4{,}760.4 & 5{,}151.3 \\
\bottomrule
\end{tabular}
\end{table}

For M1, adding ln(review\_age) raises pseudo-$R^2$ from 0.201 to 0.213
and lowers AIC; dropping ln(review\_age) from the main specification is
a conservative choice motivated by the risk of platform-era
contamination (older reviews may reflect an earlier generation of
Naver's rating interface rather than true temporal effects). For M2,
dropping season\_autumn costs about 0.016 in pseudo-$R^2$ --- the
largest single-variable impact among the alternatives --- confirming
that the autumn effect is substantive and not an artefact of including
the dummy. Adding is\_tourist to M2 marginally improves fit (pseudo-
$R^2$ 0.1924 vs 0.1904); we exclude it from the main specification
because the AIC gain is small and the predictor is already in M1,
preserving the ``variety between models'' principle motivated in \S2.2.

\sect{V. Conclusion and Implications}

\subsect{5.1 Summary of Findings}
\begin{itemize}[leftmargin=14pt, itemsep=2pt]
  \item Waiting at the entrance is the strongest and most consistent
        negative predictor in both models. It roughly halves the odds
        of a return visit and triples the odds of a LOW rating relative
        to a TOP rating.
  \item Reviewer-identity variables (follower count, prior review count)
        are substantial positive predictors of revisit; they are weak or
        null in the rating model.
  \item Tourist-purpose visits and solo visits enter revisit with
        opposite signs: tourists are much less likely to return, solo
        visitors much more.
  \item For rating, sentiment dominates (though partly mechanically),
        and autumn visits are substantially more likely to receive a
        TOP rating.
  \item The two outcomes correlate only weakly ($r = 0.08$),
        confirming that although both relate to satisfaction, they
        capture distinct aspects of it.
\end{itemize}

\subsect{5.2 Implications}
For attraction managers, the single actionable finding from this report
is that reducing wait times at the entrance should raise both the
probability of return visits and the probability of top ratings. For
platforms, the finding that reviewer-identity variables predict revisit
suggests that investment in reviewer retention — recognising active
reviewers, surfacing their histories — may have outsized returns compared
with content-side nudges aimed at individual reviews.

\subsect{5.3 Limitations and Next Steps}
\begin{itemize}[leftmargin=14pt, itemsep=2pt]
  \item \textbf{Associations only.} The design is cross-sectional and
        does not identify causal effects. The revisit indicator in
        particular is determined before the review is written, so M1
        profiles return visitors rather than explaining why anyone
        returns.
  \item \textbf{Partially mechanical sentiment--rating association.}
        Both sentiment and rating are produced in the same reviewing
        session, so the strong sentiment coefficient in M2 should be
        read as a text-level control, not a causal effect.
  \item \textbf{Imbalanced rating classes.} The MID category is small
        (18\%) and surrounded by TOP (75\%) and LOW (7\%); the
        multinomial model handles this by routing most observations to
        the TOP or LOW classes.
  \item \textbf{Place fixed effects only.} We do not include month
        fixed effects or a Hausman test comparing fixed and random
        effects; adding these is a natural next step.
  \item \textbf{Single platform.} Nothing in this analysis separates
        Naver's reviewer population from Korean reviewers more generally.
\end{itemize}

\vspace{8pt}
\subsect{References}
\begin{itemize}[leftmargin=14pt, itemsep=3pt]
  \item Bae, E., Chang, M., Park, E., \& Kim, D. (2017). The effect of
        Hallyu on tourism in Korea. \textit{Journal of Open Innovation:
        Technology, Market, and Complexity}, 3, 22.
  \item Chang, S. (2021). Food tourism in Korea. \textit{Journal of
        Vacation Marketing}.
  \item Karabila, I., Darraz, N., EL-Ansari, A., Alami, N., \& EL
        Mallahi, M. (2024). BERT-enhanced sentiment analysis for
        personalized e-commerce recommendations. \textit{Multimedia
        Tools and Applications}, 83, 56463--56488.
  \item Ko, S., Kang, S., Kang, H., \& Lee, M. J. (2018). An exploration
        of foreign tourists' perceptions of Korean food tour: a
        factor-cluster segmentation approach. \textit{Asia Pacific
        Journal of Tourism Research}, 23(8), 833--846.
  \item Lee, M., Jeong, D., \& Kim, A. J. (2026). The impact of
        information diversity on online reviews and business
        performance. \textit{Journal of Consumer Studies}, 37(1).
        \url{https://doi.org/10.35736/JCS.37.1.1}
  \item Parasuraman, A., Zeithaml, V. A., \& Berry, L. L. (1988).
        SERVQUAL: A multiple-item scale for measuring consumer
        perceptions of service quality. \textit{Journal of Retailing},
        64(1), 12--40.
\end{itemize}

% ═══════════════════════════════════════════════════════════════════════
% KOREAN VERSION
% ═══════════════════════════════════════════════════════════════════════
\clearpage
\setcounter{table}{0}
\setcounter{figure}{0}

\begin{center}
{\Large\textbf{국내 관광명소 온라인 리뷰에서 방문 맥락과\\[2pt]
리뷰 결과에 관한 실증분석}}

\vspace{12pt}

{\normalsize *소리아노 에르난데스 호르헤 아단 (호서대학교 데이터사이언스학과,
대학원생, adan\_25@naver.com)}

\end{center}

\vspace{10pt}

\noindent\textbf{\large 초록}

\vspace{4pt}

\noindent
본 연구는 한국 관광지에 대한 네이버 플레이스 온라인 리뷰에서, 플랫폼이 구조화하여
수집한 방문 맥락(동행 유형, 방문 목적, 대기 여부, 예약 여부)이 리뷰 결과에 어떻게
반영되는지를 분석하고자 한다. 선행 연구는 리뷰 텍스트의 감성이나 주제 구조에
집중해 왔으며, 플랫폼이 제공하는 방문 맥락 메타데이터를 공동 예측 집합으로 활용하여
리뷰 결과를 설명하려는 시도는 상대적으로 부족하였다. 또한 많은 연구가 단일 관광지나
단일 범주에 국한되어, 다양한 유형의 관광지를 포괄하는 리뷰 수준의 비교 분석이
부족한 실정이다. 이에 본 연구는 두 가지 연구 질문을 다룬다. 첫째, 방문 맥락과
작성자 특성이 리뷰어의 선언된 만족도(별점)에 어떻게 영향을 미치는가. 둘째, 방문
맥락과 작성자 특성이 재방문 여부에 어떻게 반영되는가. 두 결과 변수는 관측 불가능한
만족도의 서로 다른 흔적으로 간주하며, 별점은 선언된 지표, 재방문 표식은 과거
경험에 기반한 행동 지표로 해석한다. 분석 대상은 방문객 규모를 기준으로 한 층화
추출로 선정된 서울·부산 등 29개 주요 관광지의 한국어 리뷰 14{,}401건이며, 별점을
세 범주(TOP $=5$\textbigstar, MID $4$--$4.5$\textbigstar, LOW $\le 3.5$\textbigstar)로
구분한 다항 로지스틱 회귀와 재방문 여부에 대한 이항 로지스틱 회귀를 활용한다.
두 모형은 방문 맥락, 작성자 특성, 장소 통제를 공통 예측 변수로 사용한다. 본
연구는 플랫폼이 구조화한 방문 맥락을 리뷰 수준의 공동 예측 집합으로 활용하여,
선언된 만족도와 행동적 만족도의 흔적을 동시에 설명하는 분석 틀을 제시한다는
데 의의가 있다.

\vspace{6pt}

\noindent\textbf{주제어 (Keywords):} 리뷰, 방문 맥락, 재방문, 만족도,
로지스틱 회귀

\sect{I. 서론}

\subsect{1.1 연구배경}
온라인 리뷰는 여행자가 어디를 방문하고 재방문할지 결정할 때 중요한 참고
자료로 자리 잡았다. 여행 전 소비자는 TripAdvisor, Google Maps 등 주요
플랫폼과 함께 한국 시장에서는 네이버 플레이스와 카카오맵에서 리뷰, 사진,
별점을 확인한다. 이들 플랫폼은 공개하는 데이터의 종류에서 차이가 있는데,
대부분의 서구권 플랫폼은 리뷰 텍스트와 수치 별점을 공개하는 반면, 네이버
플레이스는 추가로 방문 자체에 관한 구조화된 메타데이터(동행자, 방문
목적, 대기 여부, 예약 여부)를 기록한다. 이러한 플랫폼이 수집한 필드는
각 리뷰를 기계가 판독 가능한 방문의 소규모 사례 자료로 전환시킨다.

한국의 인바운드 관광은 지난 10년간 크게 성장하였다. 팬데믹 이후 회복
과정에서도 한국은 아시아의 주요 관광 목적지 중 하나로서 다양한 체험
세분에 걸쳐 방문객을 유치하고 있다. 선행 연구는 그 성장의 일부를
한류(Bae, Chang, Park \& Kim, 2017)와 한국 음식 관광과 같은 특정
체험 세분(Chang, 2021; Ko, Kang, Kang \& Lee, 2018)에 기인하는
것으로 설명한다. 국내 리뷰 플랫폼으로 가장 활발히 사용되는 네이버
플레이스는 매년 관광지에 대한 수백만 건의 리뷰를 축적하며, 리뷰
수준의 분석을 위한 풍부한 자료원을 제공한다.

이러한 맥락에서, 방문자가 관광지를 \emph{어떻게} 평가하는지, 그리고
\emph{재방문할지 여부}를 결정짓는 맥락적 요인을 이해하는 것은 관광지
운영자에게는 실무적으로, 소비자행동 연구자에게는 학문적으로 모두 중요한
과제이다. 기존 연구들은 대부분 리뷰 텍스트를 독립된 문서로 분석하거나
업체 단위로 집계하여 리뷰 자체를 분석 단위로 삼지 않았다. 네이버가
수집하는 구조화된 방문 맥락 메타데이터는 리뷰 각각을 방문 맥락과 함께
온전하게 모형화할 수 있는 가능성을 처음으로 열어준다.

본 보고서에서는 \textbf{방문 맥락}을 플랫폼이 수집한 방문의 차원 집합
---동행 유형(혼자, 배우자/연인, 친구, 아이, 가족, 동료), 방문 목적
(여행, 나들이, 일상, 데이트, 사교), 대기 여부(이항), 예약 여부(이항)---
으로 정의한다. 이러한 차원은 다른 리뷰 생태계에서는 거의 기록되지 않는
플랫폼 구조적 메타데이터이며, 본 연구 전반의 핵심 예측 변수 집합을
구성한다. 여기에 작성자 측 변수(팔로워 수, 과거 리뷰 수, 리뷰에 첨부된
사진 수)와 텍스트 기반 감성 점수를 더하여, 각 리뷰를 방문·작성자·텍스트
특징의 간결한 집합으로 표현한다.

\subsect{1.2 연구 질문}
본 연구는 다음 두 가지 연구 질문을 다룬다.
\begin{enumerate}[leftmargin=16pt, itemsep=2pt]
  \item 방문 맥락과 작성자 특성이 리뷰어가 재방문 방문자일 확률에
        어떻게 영향을 미치는가?
  \item 방문 맥락과 작성자 특성이 각 별점 범주(5점, 4--4.5점,
        3.5점 이하)의 확률에 어떻게 영향을 미치는가?
\end{enumerate}

두 결과 변수는 잠재적 만족도 개념의 서로 다른 관측 흔적으로 해석된다.
별점은 작성 시점에 \emph{선언된} 만족도를 포착하는 반면, 재방문 표식은
이전 방문들에 걸쳐 누적된 만족도의 \emph{행동적} 흔적을 포착한다. 거의
동일한 예측 변수 집합을 사용하여 두 결과를 병렬적으로 모형화함으로써,
두 모형 간 계수의 차이가 모형 설정의 차이가 아닌 결과 변수 자체의
차이에서 비롯되도록 설계하였다.

\subsect{1.3 연구 기여}
본 연구는 세 가지 기여를 한다. \emph{(i)~플랫폼 구조적 방문 맥락의
공동 예측 집합 활용.} 네이버 플레이스의 구조화된 메타데이터(동행, 목적,
대기, 예약)를 리뷰 수준 예측 변수로 활용하여, 리뷰 텍스트만을 분석에
한정하지 않는다. \emph{(ii)~두 결과의 병렬 모형화.} 공유된 예측 변수
집합을 바탕으로 두 개의 회귀를 추정하여, 각 맥락적 요인이 재방문과
별점에 미치는 영향을 직접적으로 비교할 수 있도록 한다. \emph{(iii)~
3범주 별점에 대한 다항 로지스틱 회귀.} 별점에 대한 순서형 로지스틱
설정이 비례오즈 가정에 대한 Brant 검정에서 기각되었기 때문(\S3.5
참조), 범주별로 계수를 자유롭게 추정할 수 있는 다항 로지스틱 회귀를
채택하였다. 이를 통해 주요 예측 변수가 별점 범주별로 비대칭적 효과를
가진다는 사실을 드러낸다.

\sect{II. 이론적 배경}

\subsect{2.1 방문 맥락과 리뷰 행동}
이민영·정동영·김지영(2026)은 네이버 스마트플레이스 자료를 업체-주
단위로 집계한 자료($N = 2{,}662{,}298$)를 활용하여, 리뷰 측 특성이
업체 수준 성과에 영향을 미친다는 점을 보여주었다. 이들의 분석은 정보
다양성(리뷰 전반에 걸친 키워드 내용의 고유성)과 리뷰 감성이 모두
플랫폼 트래픽(페이지뷰, 총 클릭 수)을 예측함을 확인하고, 2023년 3월
네이버의 키워드 확장 업데이트와의 중요한 상호작용을 보고한다. 그들의
기여는 리뷰의 양뿐 아니라 \emph{내용}이 중요하다는 점을 입증한 데
있다. 다만 업체-주 수준으로의 집계는 리뷰 자체와 리뷰가 작성된 방문
맥락을 직접 관측하지 않는다는 한계를 가진다. 본 연구는 그들의 통찰을
리뷰 수준으로 확장하여, 리뷰 수준의 예측 변수가 두 행동적 결과
(재방문과 별점)를 어떻게 설명하는지 묻는다.

관광 세분화 연구는 방문 목적과 동행 유형을 방문자의 경험과 평가에
가장 강한 영향을 미치는 두 차원으로 다룬다. Ko 외(2018)는 요인-군집
세분화를 통해 외국인 관광객의 한식 투어 인식을 분류하였고, Chang
(2021)은 음식 관광 자체를 독자적인 체험 세분으로 다루었으며, Bae
외(2017)는 관광 세분 성장을 한류와 연결한다. 본 연구는 이러한 선행
연구를 토대로 방문 목적(``관광'' 지시 변수)과 동행 유형(``혼자''
지시 변수)을 리뷰 수준 예측 변수에 포함하였다.

대기 경험과 예약 여부는 서비스 접점의 운영적 신호로서, 서비스 품질
문헌(SERVQUAL; Parasuraman, Zeithaml \& Berry, 1988)에서 중요하게
다루어진다. SERVQUAL은 서비스 품질의 다섯 차원(유형성, 신뢰성, 반응성,
확신성, 공감성)을 식별하는데, 입구 대기는 반응성의 직접적인 위반이며,
관광·호텔 분야 선행 연구는 대기가 만족도와 행동 의도에 종종 큰 영향을
미친다는 점을 보여준다. 반대로 예약은 계획성과 몰입의 신호이며 서비스에
대한 기본 기대 수준을 높이는 요인으로 작용한다.

\subsect{2.2 잠재 구성개념으로서의 만족도}
소비자행동 문헌에서 고객 만족도는 복수의 관측 가능한 지표를 지닌
\emph{잠재적} 구성개념으로 흔히 개념화된다. 고전적 측정 관점(예: Oliver
의 기대-불일치 모형)은 만족도를 별점, 텍스트 감성, 재구매 의도, 재방문
등 다수의 행동적·선언적 흔적으로부터 추론되는 관측 불가능한 심리
상태로 간주한다. 단일 흔적만으로는 잠재 개념을 온전히 드러낼 수 없으며,
흔적들은 일반적으로 중간 정도의 상관관계만을 가진다.

리뷰 플랫폼 환경에서는 두 가지 특히 유용한 관측 지표가 데이터에
공존한다. 별점(작성 시점에 생성되는 \emph{선언적} 지표)과 재방문 표식
(과거 방문을 집계하는 \emph{행동적} 지표)이다. 두 지표는 관련되어
있으나 반드시 함께 움직일 필요는 없다. 리뷰어는 다섯 별을 부여하면서도
재방문하지 않을 수 있고(예: 당일치기 여행자가 그 가치는 인정하지만
돌아갈 이유가 없는 경우), 혹은 최고점이 아닌 평가를 하면서도 반복적으로
방문할 수 있다(예: 지역 주민이 열정적이지는 않지만 유용하게 활용하는
경우). 본 연구는 두 지표를 거의 동일한 예측 변수 집합으로 병렬적으로
모형화하여, 두 모형 간 계수의 차이가 모형 설정이 아닌 결과 변수 자체의
차이에 기인하도록 설계하였다.

본 표본에서 두 결과 변수 간의 약한 경험적 상관($r = 0.08$,
\S3.4 참조)은 이러한 프레이밍을 뒷받침한다. 재방문과 별점이 모두
만족도를 반영하기는 하나, 서로 다른 측면을 포착하며 공유 분산은 1\%
미만이다. 따라서 두 모형 간 계수 패턴의 발산은 모형화 선택의 결과가
아니라, 서로 다른 맥락적 요인이 두 관측 흔적으로 어떻게 경로화되는지에
관한 실질적 발견이다.

\subsect{2.3 감성 분류와 리뷰 유용성}
리뷰 텍스트 감성은 BERT 계열 한국어 감성 분류기
(\texttt{hun3359/klue-bert-base-sentiment})로 산출하였다. 이 모델은
한국어 소비자 리뷰 데이터로 파인튜닝된 모델이다. BERT 기반 감성 모델은
전통적인 bag-of-words나 규칙 기반 접근 대비 소비자 리뷰 환경에서 크게
우수한 성능을 보이며, 제품 추천 환경에서는 91\%의 정확도를 달성한
사례가 있다(Karabila 외, 2024). 본 연구는 감성 점수를 텍스트 기반
예측 변수로 다루며, 두 결과 모형 모두에 포함한다.

별점 결과에서는 감성과 별점 간의 연관이 부분적으로 기계적임을 명시한다.
두 변수 모두 동일 리뷰어가 동일 세션에서 생성한 것이므로, 긍정적으로
읽히는 리뷰와 높은 별점은 공통된 저작 과정을 공유한다. 따라서 별점
모형의 감성 계수는 텍스트 수준 통제 변수로 해석해야 하며, 인과적
효과로 해석해서는 안 된다. 재방문 결과에서는 기계적 중복이 더 약하다
---재방문 표식은 이번 방문뿐 아니라 여러 방문의 경험을 반영한다
---따라서 재방문 모형의 감성 계수는 리뷰 내용 효과로 더 해석 가능하다.

\subsect{2.4 연구 공백}
기존 연구들은 리뷰 텍스트를 독립적으로 분석하거나 업체-주 단위 성과로
집계하는 데 머물러 있으며, 플랫폼이 수집한 방문 맥락을 리뷰 수준 결과와
연결한 연구는 많지 않다. 또한 별점을 순서형 결과로 모형화할 때, 순서형
로지스틱의 비례오즈 가정은 명시적으로 검정되지 않는 경우가 많아 범주별
효과의 중요한 차이를 가릴 수 있다. 본 연구는 두 공백을 모두 다룬다.
네이버가 수집한 방문 맥락을 리뷰 수준 예측 집합으로 사용하고, 별점은
비례오즈 가정을 명시적으로 검정한 뒤 다항 결과로 모형화한다(\S3.5
참조). 이 조합은 거의 동일한 예측 변수를 사용하면서도 세 범주에 걸쳐
이질적 효과를 포착할 수 있게 한다.

\sect{III. 연구방법}

\subsect{3.1 자료}
서울과 부산을 포함한 한국의 29개 주요 관광지에 대해, 각 관광지별로
개별적으로 네이버 플레이스 리뷰를 수집하였다. 총 14{,}401건의 한국어
리뷰를 분석에 사용하며, 그중 11{,}165건은 리뷰 텍스트를 포함하고
4{,}044건은 텍스트와 별점을 모두 포함한다. 후자가 별점 모형(M2)의
표본이다. 관광지는 유형별 범주(궁궐, 전통마을, 시장, 공원, 박물관,
해변, 테마파크 등)에 걸친 층화 추출로 선정하였으며, 각 관광지는 최소
30건의 한국어 리뷰를 포함한다.

\subsect{3.2 변수}
각 리뷰는 하나의 관측치가 된다. 종속 변수는 \textbf{재방문}
(리뷰어가 2회 이상 방문한 경우 1인 이항 지시 변수)과 \textbf{별점}
(원본 0.5--5.0점 척도를 세 범주로 재분류)이다. 별점 범주는 다음과 같다.
\begin{itemize}[leftmargin=14pt, itemsep=0pt]
  \item \textbf{LOW}: 별점 $\le 3.5$\textbigstar (276건, 6.8\%)
  \item \textbf{MID}: 별점 $\in \{4.0, 4.5\}$\textbigstar (728건, 18.0\%)
  \item \textbf{TOP}: 별점 $= 5.0$\textbigstar (3{,}040건, 75.2\%)
\end{itemize}

예측 변수는 세 범주의 플랫폼 수집 필드에서 도출된다. 즉 \emph{방문
맥락}(리뷰어가 보고한 방문 관련 필드), \emph{작성자 특성}(리뷰어
프로파일의 플랫폼 수준 속성), 그리고 파생된 텍스트 기반 예측 변수
(\emph{감성})이다. 표~\ref{tab:vars-ko}는 각 모형에서 사용되는 예측
변수를 열거한다.

\begin{table}[H]\centering
\caption{모형별 예측 변수. 세 개 변수(감성, 대기, ln(팔로워))는
두 모형 모두에 공유되며, 나머지는 하나의 모형에만 포함된다.}
\label{tab:vars-ko}
\footnotesize
\begin{tabular}{@{}p{3.4cm}p{3.6cm}p{3.4cm}cc@{}}
\toprule
변수 & 설명 & 척도 & M1 & M2 \\
\midrule
감성 (sentiment)                & 텍스트의 BERT 감성          & 연속 $[-1, 1]$       & \checkmark & \checkmark \\
대기 (waited)                   & 입구에서 대기함            & 이항 $\{0, 1\}$       & \checkmark & \checkmark \\
ln(1+팔로워)                     & 리뷰어의 팔로워 수         & 연속                   & \checkmark & \checkmark \\
\midrule
관광 (is\_tourist)               & 방문 목적이 ``여행''        & 이항 $\{0, 1\}$       & \checkmark & -- \\
혼자 (comp\_alone)               & 혼자 방문                 & 이항 $\{0, 1\}$       & \checkmark & -- \\
ln(1+리뷰 수)                    & 리뷰어의 과거 리뷰 수        & 연속                   & \checkmark & -- \\
\midrule
ln(1+사진 수)                    & 리뷰에 첨부된 사진 수         & 연속                   & -- & \checkmark \\
가을 (season\_autumn)            & 방문이 가을인 경우            & 이항 $\{0, 1\}$       & -- & \checkmark \\
\bottomrule
\end{tabular}
\end{table}

변수 선택은 두 결과 변수에 대해 80개 이상의 사양을 비교한 완전탐색을
통해 이루어졌다. 최종 사양은 각 결과별로 통계적으로 유의하며 포함
시 AIC가 개선되는 변수들을 유지한다. 세 개의 공유 예측 변수는 후보
풀의 모든 부분집합에서 두 결과 변수 모두에 유의한 것으로 나타난
유일한 변수들이다.

\subsect{3.3 기술 통계}
표~\ref{tab:desc-ko}는 종속 변수와 예측 변수의 기술 통계를 보여준다.

\begin{table}[H]\centering
\caption{종속 변수 및 예측 변수의 기술 통계.}
\label{tab:desc-ko}
\footnotesize
\begin{tabular}{@{}lrrrrrr@{}}
\toprule
변수                   & N (M1)   & 평균   & 표준편차 & 최소값  & 최대값 & N (M2) \\
\midrule
재방문 (0/1)           & 11,165 & 0.226 & 0.418 & 0 & 1 & --- \\
5점 (0/1)              & ---    & 0.752 & 0.432 & 0 & 1 & 4,044 \\
LOW ($\le 3.5$\textbigstar)  & --- & 0.068 & --- & 0 & 1 & 276 \\
MID (4--4.5\textbigstar)    & --- & 0.180 & --- & 0 & 1 & 728 \\
TOP (5\textbigstar)         & --- & 0.752 & --- & 0 & 1 & 3,040 \\
대기                   & 11,165 & 0.079 & 0.270 & 0 & 1 & 4,044 \\
관광                   & 11,165 & 0.163 & 0.369 & 0 & 1 & 4,044 \\
혼자                   & 11,165 & 0.110 & 0.312 & 0 & 1 & 4,044 \\
가을                   & 11,165 & 0.301 & 0.459 & 0 & 1 & 4,044 \\
ln(1+리뷰)             & 11,165 & 5.386 & 1.927 & 0.69 & 9.86 & 4,044 \\
ln(1+팔로워)           & 11,165 & 1.895 & 2.245 & 0.00 & 9.17 & 4,044 \\
ln(1+사진)             & 11,165 & 0.100 & 0.375 & 0.00 & 3.05 & 4,044 \\
감성                   & 11,165 & 0.653 & 0.457 & $-0.97$ & 0.99 & 4,044 \\
\bottomrule
\end{tabular}
\end{table}

전체 리뷰 중 약 23\%가 재방문 리뷰이며, 별점 리뷰의 75\%가 5점을
받았다. 입구에서 대기한 비율은 8\%에 불과하며, 관광 목적 방문이
16\%, 혼자 방문이 11\%이다. 가을 방문이 전체의 30\%를 차지하며,
감성은 긍정 쪽으로 치우쳐 있다(평균 $0.65$).

\subsect{3.4 상관 분석 및 다중공선성}
8개 예측 변수 간 쌍별 상관은 ln(리뷰)와 ln(팔로워) 간의
$r = 0.73$을 제외하고 모두 절댓값 0.40 미만이다. 분산팽창계수(VIF)는
모두 2.5 미만으로, 다중공선성 문제가 없음을 확인하였다. 가장 높은
상관이 나타난 두 변수(ln(리뷰), ln(팔로워))는 M1에만 포함되고 M2에는
단독으로 포함되지 않으므로, 두 모형 모두에서 공선성 문제는 발생하지
않는다.

본 연구의 핵심 프레이밍에 중요한 결과로, 두 종속 변수 자체는 약한
상관만을 보인다. 공유 표본($N = 4{,}044$)에서 재방문과 5점 지시
변수 간 상관은 $r = 0.08$ ($p < 0.001$)이다. 이 낮은 상관은 두 결과를
별도로 모형화하는 경험적 근거이다. 두 변수 모두 잠재적 만족도 구성
개념의 지표로 간주되지만, 서로 다른 측면을 포착하며 공유 분산은
약 0.6\%에 불과하다.

\subsect{3.5 모형}
두 모형 모두 관광지 고정효과 $\delta_j$(28개 더미, 경복궁 참조)와
작성자(\texttt{user\_id}) 수준 군집 표준오차를 포함한다. $i$는 리뷰,
$j$는 장소의 인덱스이며, $\Lambda(z) = 1/(1+e^{-z})$는 로지스틱 함수이다.

\textbf{모형 1 --- 이항 로지스틱 회귀 (재방문).} 종속 변수
$\text{Revisit}_{ij}$는 리뷰어의 방문 횟수가 2회 이상이면 $1$, 그렇지
않으면 $0$이다.

\begin{align*}
\operatorname{logit} P(\text{Revisit}_{ij} = 1)
   = \; &\alpha
       + \beta_{1}\,\text{waited}_{ij}
       + \beta_{2}\,\text{is\_tourist}_{ij}
       + \beta_{3}\,\text{sentiment}_{ij} \\
    &+ \beta_{4}\,\text{comp\_alone}_{ij}
       + \beta_{5}\,\ln(1 + \text{reviews}_{ij})
       + \beta_{6}\,\ln(1 + \text{followers}_{ij})
       + \delta_{j}
\end{align*}

\textbf{모형 2 --- 다항 로지스틱 회귀 (별점).} $Y_{ij} \in
\{\text{LOW}, \text{MID}, \text{TOP}\}$를 3범주 별점이라 하고,
TOP를 참조 범주로 설정한다. 각 비참조 범주 $k \in \{\text{MID},
\text{LOW}\}$에 대해:

\begin{align*}
\log\!\frac{P(Y_{ij} = k)}{P(Y_{ij} = \text{TOP})}
   = \; &\alpha_{k}
       + \gamma_{k1}\,\text{sentiment}_{ij}
       + \gamma_{k2}\,\text{waited}_{ij}
       + \gamma_{k3}\,\ln(1 + \text{followers}_{ij}) \\
    &+ \gamma_{k4}\,\ln(1 + \text{photos}_{ij})
       + \gamma_{k5}\,\text{season\_autumn}_{ij}
       + \delta_{jk}
\end{align*}

각 비참조 범주는 고유의 절편 $\alpha_k$, 고유의 기울기 계수 벡터
$(\gamma_{k1}, \ldots, \gamma_{k5})$, 그리고 고유의 장소 고정효과
$\delta_{jk}$를 가진다.

\textbf{왜 다항 로지스틱인가?} 초기 사양에서 M2를 순서형 로지스틱
회귀로 추정하였으나, Brant 검정에서 전체 모형의 비례오즈 가정이
기각되었다($\chi^2 = 18.8$, df$=7$, $p = 0.009$). 특히 감성과
ln(review\_age)이 개별적으로 비례오즈를 위반하였다($p < 0.01$).
다항 로지스틱 회귀는 각 범주 비교에 대해 별도의 계수를 추정함으로써
이 가정을 제거한다. 다만 트레이드오프로, 별점의 순서적 구조는 적합에
강제되지 않는다.

\sect{IV. 연구결과}

\subsect{4.0 모형 적합도 요약}
표~\ref{tab:fit-ko}는 두 모형의 적합도를 나란히 요약한다. 두 모형
모두 소비자 연구에서 이산 선택 모형의 ``좋음'' 범위에 도달한다
(pseudo-$R^2$ 0.10 이상, McFadden에 따르면 0.20--0.40이 ``매우 좋음''
수준).

\begin{table}[H]\centering
\caption{모형 적합도 요약.}
\label{tab:fit-ko}
\footnotesize
\begin{tabular}{@{}l l l@{}}
\toprule
통계량                   & M1 재방문             & M2 별점 \\
\midrule
모형 유형                 & 이항 로지스틱            & 다항 로지스틱 \\
결과 변수                 & 재방문 (0/1)          & 별점 (LOW / MID / TOP) \\
참조 범주                 & 재방문 아님 (0)         & TOP (5\textbigstar) \\
관측치 수                 & 11,165               & 4,044 \\
예측 변수 수              & 6                    & 5 (범주당) \\
장소 고정효과             & 28                   & 28 (범주당) \\
pseudo-$R^2$              & 0.201                & 0.190 \\
AIC                       & 9{,}601.5            & 4{,}753.6 \\
BIC                       & 9{,}857.7            & 5{,}157.1 \\
AUC / 정확도 (CV)         & 0.80                 & 0.77 \\
Within-$\pm 1$ 정확도 (M2)& ---                  & 0.95 \\
\bottomrule
\end{tabular}
\end{table}

\subsect{4.1 모형 1 --- 재방문}
표~\ref{tab:m1-ko}는 모형 1의 계수 추정 결과를 보여준다. 6개 예측
변수 모두 $p < 0.01$에서 유의하다. pseudo-$R^2$는 $0.201$이며,
작성자 기준 5-fold 교차검증 결과 AUC는 0.78이다.

\begin{table}[H]\centering
\caption{모형 1: 재방문의 이항 로지스틱 회귀. $N = 11{,}165$;
pseudo-$R^2 = 0.201$. 장소 고정효과 포함; 작성자 수준 군집 표준오차.}
\label{tab:m1-ko}
\footnotesize
\begin{tabular}{@{}lrrr@{}}
\toprule
변수 & OR & 95\% CI & $p$ \\
\midrule
대기 (waited)                 & 0.54 & [0.37, 0.79] & 0.001 ** \\
관광 (is\_tourist)             & 0.55 & [0.42, 0.73] & $<\,.001$ *** \\
감성 (sentiment)               & 0.74 & [0.64, 0.87] & $<\,.001$ *** \\
혼자 (comp\_alone)             & 2.26 & [1.67, 3.07] & $<\,.001$ *** \\
ln(1+리뷰)                     & 1.31 & [1.17, 1.48] & $<\,.001$ *** \\
ln(1+팔로워)                   & 1.23 & [1.13, 1.35] & $<\,.001$ *** \\
\bottomrule
\end{tabular}
\end{table}

가장 강력한 부정적 효과는 \textbf{입구 대기}이다. 대기한 리뷰어의
리뷰는 재방문 리뷰일 확률이 약 절반 수준이다(OR $= 0.54$). 이는
본 보고서의 가장 명확한 운영적 발견이며, 관광지 운영자가 직접
조치할 수 있는 유일한 변수이다.

두 \textbf{작성자 정체성 변수}가 양의 방향으로 나타난다.
ln(1$+$리뷰) (OR $= 1.31$)와 ln(1$+$팔로워) (OR $= 1.23$)이다. 플랫폼
활동 이력이 많거나 네트워크가 큰 리뷰어일수록 재방문 리뷰일 가능성이
높다. 이는 자기선택 효과로 해석할 수 있다. 적극적인 리뷰어는 온라인
활동이 많고 재방문도 많이 하는 경향이 있다.

\textbf{방문 유형 효과}는 크고 반대 방향이다. 관광 목적 방문은
재방문일 확률이 45\% 낮으며(OR $= 0.55$), 혼자 방문은 재방문 확률을
두 배 이상 높인다(OR $= 2.26$). 혼자 방문자가 지역 주민이거나 자주
돌아오는 애호가일 수 있음을 시사한다.

감성은 M1에서 \emph{음의} 방향으로 들어온다(OR $= 0.74$). 더 긍정적인
감성은 재방문 리뷰일 확률을 약간 낮춘다. 신규 방문자는 비교 대상이
없기 때문에 더 균일하게 긍정적인 리뷰를 쓰는 반면, 재방문자는 더
세밀한 평가를 담는다는 해석이 가능하다.

\subsect{4.2 모형 2 --- 별점}
표~\ref{tab:m2-ko}는 모형 2 계수를 두 열의 오즈비로 보여준다
(MID 대 TOP, LOW 대 TOP). pseudo-$R^2$는 0.190이다.

\begin{table}[H]\centering
\caption{모형 2: 별점 범주의 다항 로지스틱 회귀. $N = 4{,}044$;
pseudo-$R^2 = 0.190$. TOP(5점)을 참조 범주로 하는 오즈비. 장소
고정효과 포함; 작성자 수준 군집 표준오차.}
\label{tab:m2-ko}
\footnotesize
\begin{tabular}{@{}l rr rr@{}}
\toprule
                            & \multicolumn{2}{c}{MID 대 TOP} & \multicolumn{2}{c}{LOW 대 TOP} \\
\cmidrule(lr){2-3} \cmidrule(l){4-5}
변수                        & OR    & $p$          & OR    & $p$ \\
\midrule
감성                        & 0.44  & $<\,.001$ ***  & 0.04  & $<\,.001$ *** \\
대기                        & 1.21  & 0.326        & 3.10  & $<\,.001$ *** \\
ln(1+팔로워)                & 0.95  & 0.287        & 0.71  & $<\,.001$ *** \\
ln(1+사진)                  & 0.81  & 0.197        & 0.44  & 0.043 *      \\
가을                        & 0.39  & $<\,.001$ ***  & 0.38  & $<\,.001$ *** \\
\bottomrule
\end{tabular}
\end{table}

다항 분해는 예측 변수가 두 범주 비교에서 \textbf{비대칭적으로 작용함}을
드러낸다. \textbf{감성}은 LOW 대 TOP(OR $= 0.04$)에서 MID 대 TOP
(OR $= 0.44$)보다 훨씬 강한 효과를 보인다. 감성과 별점이 동일 세션에서
생성되므로 이 연관은 부분적으로 기계적이며, 인과적 효과로 해석해서는
안 된다.

\textbf{대기}는 두드러진 비대칭 패턴을 보인다. MID 대 TOP와는
유의하지 않지만($p = 0.33$), LOW 대 TOP의 확률을 크게 높인다
(OR $= 3.10$, $p < 0.001$). 실무적으로 대기는 ``평범함 신호''가
아닌 ``실패 신호''이다. 대기는 매우 낮은 별점의 강력한 예측 변수이지만,
중간 수준 별점으로 리뷰어를 밀어내지는 않는다.

\textbf{가을 효과.} 가을(9월--11월) 방문은 MID나 LOW 범주일 확률이
크게 낮다(OR $\approx 0.4$). 이는 한국의 가을이 단풍과 쾌적한 날씨로
본 표본의 관광지 범주에 피크 시즌으로 인식되는 것과 일치한다.

\subsect{4.3 두 모형의 통합적 해석}
두 모형을 나란히 비교하면 수렴과 발산이 모두 드러난다. 공유 예측 변수
(대기, 감성, ln(팔로워))는 \S2.2에서 제안한 \emph{발산 흔적} 해석을
뒷받침한다. 대기는 재방문에 강하게 부정적이며 별점에서는 ``실패 신호''이다.
감성은 재방문에 약하게 부정적이지만 높은 별점에 거의 기계적으로
긍정적이다. 팔로워 수는 양쪽 모두에 양의 방향이며, 방문 특성이 아닌
리뷰어 정체성을 반영한다.

두 결과 간의 낮은 상관($r = 0.08$)은 공유 예측 변수가 두 모형에
동일한 서사를 강제하지 않음을 확인시켜 준다. 두 모형은 \S2.2에서
서술한 만족도의 선언적 흔적과 행동적 흔적의 구별을 반영하는 서로
다른 계수 패턴을 복원한다.

\subsect{4.4 교차검증 세부 결과}
표~\ref{tab:cv-ko}는 두 모형의 fold별 교차검증 결과를 보여준다. M1의
AUC는 fold에 걸쳐 0.75에서 0.84까지 변동하며(평균 0.78), M2의 정확도는
0.76에서 0.80까지 변동한다(평균 0.77). 어느 fold도 눈에 띄게 저하된
성능을 보이지 않아, 모형이 리뷰어 하위 집합 전반에 걸쳐 일관되게
일반화됨을 시사한다.

\begin{table}[H]\centering
\caption{교차검증 세부 결과 (작성자 기준 5-fold).}
\label{tab:cv-ko}
\footnotesize
\begin{tabular}{@{}c rrrrrr | rr@{}}
\toprule
     & \multicolumn{6}{c|}{M1 재방문} & \multicolumn{2}{c}{M2 별점} \\
Fold & N test  & AUC   & Brier & 정확도 & 균형 정확도 & F1    & 정확 & $\pm 1$ \\
\midrule
1    & 2{,}233 & 0.839 & 0.136 & 0.777 & 0.574     & 0.281 & 0.755 & 0.953 \\
2    & 2{,}233 & 0.797 & 0.131 & 0.824 & 0.659     & 0.477 & 0.802 & 0.962 \\
3    & 2{,}233 & 0.751 & 0.153 & 0.781 & 0.591     & 0.335 & 0.759 & 0.942 \\
4    & 2{,}233 & 0.759 & 0.147 & 0.785 & 0.573     & 0.284 & 0.768 & 0.958 \\
5    & 2{,}233 & 0.756 & 0.145 & 0.806 & 0.629     & 0.417 & 0.757 & 0.954 \\
\midrule
평균 & 11{,}165 & 0.780 & 0.142 & 0.795 & 0.605 & 0.359 & 0.768 & 0.954 \\
\bottomrule
\end{tabular}
\end{table}

\subsect{4.5 강건성 검증}
표~\ref{tab:rob-ko}는 두 모형의 대안적 사양을 보고한다. 주 모형의
적합도 통계량이 예측 집합의 작은 변화에 안정적임을 확인한다.

\begin{table}[H]\centering
\caption{강건성 대안. 각 행은 하나의 사양을 추정하며, 각 패널의 첫
행은 \S4.1--4.2에서 보고한 주 모형이다.}
\label{tab:rob-ko}
\footnotesize
\begin{tabular}{@{}l ccrr@{}}
\toprule
사양                                   & 변수 수 & pseudo-$R^2$ & AIC     & BIC     \\
\midrule
\multicolumn{5}{@{}l}{\textit{M1 재방문 대안}} \\
주 모형 (6)                             & 6      & 0.2010       & 9{,}601.5 & 9{,}857.7 \\
+ ln(review\_age)                      & 7      & 0.2133       & 9{,}457.6 & 9{,}721.2 \\
+ season\_autumn                        & 7      & 0.2013       & 9{,}600.5 & 9{,}864.0 \\
$-$ comp\_alone                         & 5      & 0.1912       & 9{,}717.5 & 9{,}966.4 \\
$-$ sentiment                           & 5      & 0.1989       & 9{,}624.8 & 9{,}873.7 \\
\midrule
\multicolumn{5}{@{}l}{\textit{M2 별점 대안}} \\
주 모형 (5)                             & 5      & 0.1904       & 4{,}753.6 & 5{,}157.1 \\
+ ln(review\_age)                      & 6      & 0.1920       & 4{,}748.5 & 5{,}164.6 \\
+ is\_tourist                           & 6      & 0.1924       & 4{,}746.1 & 5{,}162.2 \\
+ comp\_alone                          & 6      & 0.1904       & 4{,}757.6 & 5{,}173.7 \\
$-$ season\_autumn                     & 4      & 0.1749       & 4{,}838.5 & 5{,}229.4 \\
$-$ ln\_photos                         & 4      & 0.1885       & 4{,}760.4 & 5{,}151.3 \\
\bottomrule
\end{tabular}
\end{table}

M1의 경우 ln(review\_age)을 추가하면 pseudo-$R^2$가 0.201에서 0.213으로
상승하지만, 플랫폼 세대 혼란(오래된 리뷰가 네이버 평점 인터페이스의
이전 세대를 반영할 위험)을 고려하여 보수적으로 주 사양에서 제외하였다.
M2의 경우 season\_autumn을 제거하면 pseudo-$R^2$가 약 0.016 하락하며,
이는 가을 효과가 실질적이며 더미 포함의 인위적 결과가 아님을 확인시킨다.

\sect{V. 결론 및 시사점}

\subsect{5.1 연구결과 요약}
\begin{itemize}[leftmargin=14pt, itemsep=2pt]
  \item 입구 대기는 두 모형 모두에서 가장 강하고 일관된 부정적
        예측 변수이다. 재방문 확률을 약 절반으로 줄이고, LOW
        별점(TOP 대비)의 확률을 세 배 늘린다.
  \item 작성자 정체성 변수(팔로워 수, 과거 리뷰 수)는 재방문의
        실질적인 양의 예측 변수이며, 별점 모형에서는 약하거나 무효하다.
  \item 관광 목적 방문과 혼자 방문은 재방문에서 반대 방향으로 작용한다.
        관광객은 재방문할 확률이 훨씬 낮고, 혼자 방문자는 훨씬 높다.
  \item 별점의 경우 감성이 지배적이며(부분적으로 기계적임을 유의),
        가을 방문은 TOP 별점일 확률을 크게 높인다.
  \item 두 결과는 약한 상관($r = 0.08$)만을 보이며, 두 변수 모두
        만족도와 관련되지만 서로 다른 측면을 포착함을 확인시킨다.
\end{itemize}

\subsect{5.2 시사점}
관광지 운영자에게 본 보고서의 유일하게 실행 가능한 발견은 입구
대기 시간 단축이 재방문 확률과 TOP 별점 확률을 모두 높일 수 있다는
점이다. 플랫폼의 경우, 작성자 정체성 변수가 재방문을 예측한다는 사실은
적극적인 리뷰어 유지와 그 이력을 노출하는 투자가 콘텐츠 측의 넛지보다
더 큰 수익을 가져올 수 있음을 시사한다.

\subsect{5.3 한계 및 향후 과제}
\begin{itemize}[leftmargin=14pt, itemsep=2pt]
  \item \textbf{연관 분석에 한정.} 본 연구는 횡단면적 설계이며
        인과적 효과를 식별하지 않는다. 특히 재방문 지시 변수는 리뷰
        작성 이전에 이미 결정되므로, M1은 왜 사람들이 돌아오는지를
        설명하기보다 재방문자를 프로파일링한다.
  \item \textbf{감성과 별점 간의 부분적 기계적 연관.} 감성과 별점은
        동일한 리뷰 세션에서 생성되므로, M2의 강한 감성 계수는 텍스트
        수준 통제 변수로 읽어야 한다.
  \item \textbf{별점 범주의 불균형.} MID 범주는 작으며(18\%) TOP(75\%)과
        LOW(7\%) 사이에 끼어 있다. 다항 모형은 이를 TOP 또는 LOW
        범주로 대부분의 관측치를 라우팅하는 방식으로 처리한다.
  \item \textbf{장소 고정효과만 포함.} 월 고정효과나 고정 대 랜덤 효과를
        비교하는 Hausman 검정을 포함하지 않았다. 이를 추가하는 것이
        자연스러운 다음 단계이다.
  \item \textbf{단일 플랫폼.} 본 분석에서 네이버의 리뷰어 모집단을
        일반 한국 리뷰어와 분리하지 않았다.
\end{itemize}

\vspace{8pt}
\subsect{참고문헌}
\begin{itemize}[leftmargin=14pt, itemsep=3pt]
  \item Bae, E., Chang, M., Park, E., \& Kim, D. (2017). The effect of
        Hallyu on tourism in Korea. \textit{Journal of Open Innovation:
        Technology, Market, and Complexity}, 3, 22.
  \item Chang, S. (2021). Food tourism in Korea. \textit{Journal of
        Vacation Marketing}.
  \item Karabila, I., Darraz, N., EL-Ansari, A., Alami, N., \& EL
        Mallahi, M. (2024). BERT-enhanced sentiment analysis for
        personalized e-commerce recommendations. \textit{Multimedia
        Tools and Applications}, 83, 56463--56488.
  \item Ko, S., Kang, S., Kang, H., \& Lee, M. J. (2018). An exploration
        of foreign tourists' perceptions of Korean food tour: a
        factor-cluster segmentation approach. \textit{Asia Pacific
        Journal of Tourism Research}, 23(8), 833--846.
  \item 이민영, 정동영, 김지영. (2026). 온라인 리뷰의 정보 다양성이
        비즈니스 성과에 미치는 영향. \textit{소비자학연구}, 37(1).
        \url{https://doi.org/10.35736/JCS.37.1.1}
  \item Parasuraman, A., Zeithaml, V. A., \& Berry, L. L. (1988).
        SERVQUAL: A multiple-item scale for measuring consumer
        perceptions of service quality. \textit{Journal of Retailing},
        64(1), 12--40.
\end{itemize}

\end{document}
